% Part: first-order-logic
% Chapter: introduction
% Section: syntax

\documentclass[../../../include/open-logic-section]{subfiles}

\begin{document}

\olfileid{fol}{int}{syn}

\olsection{نحو}

سب توں پہلاں سانوں ٹھیک ٹھیک طے کرنا پینا اے کہ علامتاں دیاں
کیہڑیاں لڑیاں پہلے درجے دی منطق دے !!{sentence}s شمار ہوندیاں
نیں۔ ایہ کم اسیں بعد وچ کراں گے؛ فی الحال مثالاں نال اگے
ودھدے آں۔ پہلے درجے دی منطق دیاں بنیادی اینٹاں---یعنی
لفظیات---دے دو حصے نیں۔ پہلا حصہ اوہ علامتاں نیں جنہاں
نال اسیں کسے خاص شے ول اشارہ کردے آں یا اوہدے بارے کوئی
خاص گل کہندے آں۔ کسے شے ول اشارہ کرن لئی اسیں
!!{constant}s ورتدے آں، تے اوہدی بابت گل کرن لئی
!!{predicate}s۔ مثال لئی، کسے اک شے ول اشارہ کرن
لئی اسیں $\Obj a$ نوں !!a{constant} بنا سکدے آں،
تے اوس شے بارے !!{sentence}~$\Atom{\Obj P}{\Obj
a}$ دے ذریعے کجھ کہہ سکدے آں۔ جے تہاڈے ذہن وچ
”$\Obj a$“ تے ”$\Obj P$“ دے کجھ معنی ہون،
تاں $\Atom{\Obj P}{\Obj a}$ نوں انگریزی دے جملے
وانگ پڑھ سکدے او (تے جدوں تسی پہلی واری رسمی منطق
سکھی ہووے گی، غالباً ایویں پڑھیا وی ہووے گا)۔
جدوں پہلے درجے دی منطق دے ایہوجے سادے !!{sentence}s
تہاڈے کول ہون، تاں لفظیات دے دوجے حصے، یعنی منطقی
علامتاں (رابطے تے مسوّر)، نال ہور گنجل دار جملے
بنا سکدے او۔ مثلاً اسیں
$(\Atom{\Obj P}{\Obj a} \land \Atom{\Obj Q}{\Obj b})$
یا~$\lexists[\Obj x][\Atom{\Obj P}{\Obj x}]$
ورگے اظہار بنا سکدے آں۔

ٹھیک معنوی تعریفاں تے اپنے !!{derivation} نظاماں
دے قاعدے دین لئی، جیہڑے سخت مابعد منطقی مطالعے
لئی ضروری نیں، سانوں سب توں پہلاں ایہ ٹھیک
تعریف دینا پینی اے کہ پہلے درجے دی منطق دا
!!a{sentence} کیہنوں آکھدے آں۔ بنیادی خیال سمجھنا
اوکھا نہیں: کجھ سادے !!{sentence}s اسیں صرف
!!{predicate}s تے !!{constant}s نال بنا سکدے آں،
مثلاً~$\Atom{\Obj P}{\Obj a}$۔ فیر انہاں توں
رابطیاں تے مسوّراں نال ہور گنجل دار جملے
بناندے آں۔ پر اوہ عین کیہڑے قاعدے نیں جنہاں
مطابق اسیں ایہ گنجل دار !!{sentence}s بنا سکدے
آں؟ انہاں قاعدیاں نوں دَسنا ضروری اے؛ ورنہ
”پہلے درجے دی منطق دا !!{sentence}“ ساڈی تعریف
وچ کافی ٹھیک نہیں ہووے گا۔ ایتھے کجھ مسئلے
نیں۔ پہلا، عین اوہ لڑیاں چننا جیہڑیاں
!!{sentence}s شمار ہون۔ دوجا، ایہ کم ایس طرح
کرنا کہ اسیں \emph{سارے} !!{sentence}s بارے
ریاضیاتی ثبوت دے سکّیے۔ آخر وچ سانوں
!!{sentence}s اُتے کجھ بنیادی عملاں دی وی ٹھیک
تعریف دینی پینی اے، مثلاً ”ہر $\Obj x$،
جیہڑا~$!A$ وچ آوندا اے، اوہدی تھاں~$\Obj b$
رکھو“۔
اُکّھیڑ ایہ اے کہ پہلے درجے دی منطق دے
مسوّر تے !!{variable}s ایہ کم کیویں کرنا اے،
ایہنوں بالکل صاف نہیں رہن دیندے۔ مثال لئی،
کی $\lexists[\Obj x][\Atom{\Obj P}{\Obj a}]$
نوں !!a{sentence} شمار کرنا چاہیدا اے؟
تے $\lexists[\Obj x][\lexists[\Obj x][\Atom{\Obj P}{\Obj x}]]$
دا کیہ؟ تے ”$\Obj x$ دی تھاں~$\Obj b$ رکھو،
جیہڑا $(\Atom{\Obj P}{\Obj x} \land \lexists[\Obj x][\Atom{\Obj P}{\Obj x}])$
وچ آوندا اے“ دا نتیجہ کیہ ہونا چاہیدا اے؟

\end{document}
